% Copyright (c) 2026 by C. Klukas.
% Licensed under the MIT License. See LICENSE for details.

\paglets{} revisits a programming model from the 1990s with a runtime that
did not exist then. This chapter places it: the Aglets lifecycle it keeps,
the kinds of mobility that mobile-agent systems and process migration
offer, the actor systems that solved placement differently, and
WebAssembly as the unit that moves. It ends with the Python edition of
\paglets{} and what the C++ edition changes.

\subsection{Aglets and the agent lifecycle}\label{sec:background-aglets}

Aglets, developed at IBM's Tokyo Research Laboratory, made mobile agents
Java objects \cite{lange1997aglets,lange1998aglets}. An aglet lives in an
\emph{aglet context}, the execution environment a server such as Tahiti
offers \cite{ferrari2004manual}. Other code never holds the aglet itself;
it holds an \emph{aglet proxy}, which shields the aglet and forwards
messages to it wherever it is. The programming model is a small set of
lifecycle operations with a callback for each \cite{langeOshimaAgletApi}:

\begin{itemize}
\item \emph{creation} and \emph{cloning} (\code{onCreation},
  \code{onClone});
\item \emph{dispatch} to another context and \emph{retraction} back
  from it (\code{onDispatching}, \code{onArrival} and
  \code{onReverting});
\item \emph{deactivation} to secondary storage and \emph{activation}
  from it (\code{onDeactivating}, \code{onActivation});
\item \emph{disposal} (\code{onDisposing});
\item \emph{messages}, handled by \code{handleMessage}, sent through the
  proxy synchronously, as a future or one-way.
\end{itemize}

After creation, cloning, activation or arrival the runtime calls
\code{run()}, the aglet's own thread of control. An aglet moves as a
Java-serialized object over the Agent Transfer Protocol; the thread does
not move, so \code{run()} starts again from its beginning on the new
host. IBM's own researchers had asked in 1995 whether mobile agents are a
good idea at all \cite{harrison1995mobileAgents}; Lange and Oshima later
answered with seven reasons to move computation rather than data
\cite{lange1999seven}. The design space of such systems, from code
shipping to security, is surveyed by Karnik and Tripathi
\cite{karnik1998design}. Aglets was later released as an open-source
project on SourceForge \cite{agletsSourceForge}.

\paglets{} keeps this vocabulary. A \paglet{} is created or cloned,
dispatched, deactivated and activated, disposed; it arrives; it handles
messages; and every Aglets callback except \code{onReverting} has a
counterpart among its hooks: \code{on\_created}, \code{on\_cloned},
\code{on\_dispatching}, \code{on\_arrived}, \code{on\_deactivating},
\code{on\_activated} and \code{on\_disposing}
(Section~\ref{sec:lifecycle}). \paglets{} adds \code{on\_move\_failed},
\code{on\_undelivered} and \code{on\_recovered}, reports for which Aglets
has no callback. Two things differ at this level. There is no
\code{run()}: a \paglet{} has no thread of its own, and all its code runs
in handlers that the host calls one at a time. And there is no retraction
as a named operation: a \paglet{} that should come back is dispatched
again, by itself or by its owner from the command line
(Section~\ref{sec:tools}). The proxy survives as the \emph{endpoint}, a
capability that reaches a \paglet{} wherever it is
(Section~\ref{sec:communication}).

\subsection{Weak and strong mobility}\label{sec:background-mobility}

Fuggetta, Picco and Vigna separate the state of a running program into
its code, its data (the heap and the variables) and its execution state
(the stack and the program counter), and classify mobility by what moves
\cite{fuggetta1998mobility}. \emph{Strong} mobility moves all three: the
program continues on the new host with the statement after the jump.
Agent Tcl does this by capturing the interpreter's state \cite{gray1996agentTcl}.
\emph{Weak} mobility moves code and data only; execution restarts at an
entry point. Java systems on the standard virtual machine are weak,
because the virtual machine does not expose the state of a thread;
Aglets' restart of \code{run()} is the visible consequence.
Table~\ref{tab:mobility-kinds} puts these kinds next to process migration
and the two \paglets{} editions.

\begin{table}[htbp]
\centering\small
\caption{Kinds of mobility by what moves and where execution resumes.}
\label{tab:mobility-kinds}
\begin{tabular}{p{0.19\linewidth}p{0.35\linewidth}p{0.19\linewidth}p{0.15\linewidth}}
\toprule
Kind & What moves & Resumes at & Examples \\
\midrule
Process migration & the address space, registers, kernel state (files, sockets) & the interrupted instruction & CRIU \cite{criu}, the systems in \cite{milojicic2000migration} \\
Strong mobile agents & code, data, interpreter stack and program counter & the statement after the jump & Agent Tcl \cite{gray1996agentTcl} \\
Weak mobile agents & code (or a reference to it) and serialized object fields & \code{run()} & Aglets \cite{lange1998aglets} \\
Explicit state & a declared state record; code installed on every host & a lifecycle hook & \paglets{}, Python edition \cite{klukas2026pagletsNote} \\
Memory image & the module (by hash) and the linear memory with its globals, as bytes & \code{on\_arrived}, then the next handler & \paglets{}/cpp (this paper) \\
\bottomrule
\end{tabular}
\end{table}

In these terms \paglets{}/cpp is weak mobility: a \paglet{} moves only
between two handler calls, when no code of the \paglet{} is running, and
it resumes with the next event or message. What distinguishes it is the
data that moves. It is not a record the author declared, but the whole
linear memory of the WebAssembly instance together with its exported
mutable globals, taken as bytes. The heap, every C++ object, and the
C/C++ shadow stack that the compiler keeps in linear memory are in it.
Nothing in memory is serialized, and nothing in memory can be left
behind, because there is no list of what to take; what lives outside
memory is the subject of Section~\ref{sec:movement}. The Wasm operand
stack and the native call stack hold values only during a call, and at
the moment of the move there is no call. Moving inside a handler would
need the stack as well; the Asyncify transformation
\cite{zakai2019asyncify} and the proposed stack-switching instructions
could provide it, and mid-handler moves are a research item, not a
feature (Section~\ref{sec:discussion}).

\subsection{Process migration and actors}\label{sec:background-actors}

Process migration moves a whole operating-system process: its address
space, registers and, as far as possible, its kernel state
\cite{milojicic2000migration}. CRIU does this on Linux today, as
checkpoint and restore of process trees from user space \cite{criu}. The
unit is transparent to the program, which is its strength, and bound to
one kernel, one architecture and the same installed libraries, which is
its limit. A \paglet{} is a smaller unit with an explicit boundary: the
paglet ABI (Appendix~\ref{app:abi}). What a \paglet{} holds outside its
memory, such as an endpoint to another \paglet{} or a grant for a host
resource, is a capability with a defined mobility rule; the host
re-creates it on the destination or reports it lost
(Section~\ref{sec:movement}). The image is independent of the CPU and
the operating system; the project's feasibility milestone moved images
between macOS and Linux hosts, and between arm64 hosts and an x86-64
build of the host run under Rosetta~2, without translation
(\cite{pagletsCpp}; the milestone record in the repository's planning
documents, and Section~\ref{sec:evaluation}).

The actor model gives the other half of the lineage. An actor handles one
message at a time, creates actors and sends messages to the actors it
knows \cite{hewitt1973actor}. Erlang made this an engineering practice:
isolated processes with their own heaps, message passing as the only
interaction, supervision trees, and the rule that a failed process is
restarted rather than repaired \cite{armstrong2003thesis}. Orleans'
\emph{virtual actors} (grains) remove placement from the programmer: a
grain always exists, the runtime activates it on some server when a
message arrives and deactivates it when idle; its state can be persisted
and reloaded, and in the design as published a grain is never moved as a
running object \cite{bernstein2014orleans}. Later releases added the
migration of an idle activation to another server, its in-memory state
dehydrated on one side and rehydrated on the other
\cite{orleansGrainMigration}.

A \paglet{} is an actor in the first sense: one handler at a time, and
other \paglets{} reachable only through endpoints it was given. It
differs from grains in that placement is a decision of the \paglet{} or
its owner, who name the destination in a transfer ticket (a host, a
label, a service offer or \code{any}; for the last three the host picks a
matching host), not of a scheduler that moves \paglets{} on its own, and
the state moves as a running instance's memory rather than through a
store. An inactive \paglet{} resembles a deactivated grain: an image on
disk, restored when a message or a timer arrives. Fault tolerance
follows Erlang rather than Orleans: a zombie \paglet{}
(Section~\ref{sec:zombies}) has a replica on every host of the mesh that
offers the \code{zombies} service, and the replicas run a protocol to
take over, instead of relying on a shared store. The host restarts no
\paglet{} that trapped: as in Erlang, a crash ends the process, and
recovery is the job of another one.

\subsection{WebAssembly as the unit that moves}\label{sec:background-wasm}

WebAssembly is a portable binary instruction format with a formal
semantics, designed so that a module validates and compiles fast and runs
inside a sandbox \cite{haas2017wasm}. A paglet module declares functions,
globals, tables and exactly one linear memory (the paglet ABI requires
it, exported as \code{memory}), a contiguous array of bytes that the
module addresses with 32-bit offsets; an instance is the module together
with that memory, its globals and its tables. A C++ program compiled with
the WASI toolchain keeps its heap and its shadow stack in linear memory;
the only other state that outlives a call is the stack pointer, a mutable
Wasm global (the single mutable global of a C/C++ guest). This is the
fact on which \paglets{}/cpp rests: between two calls into the instance,
the linear memory and the mutable globals are the complete state, and the
specification defines both without reference to the host.

Runtimes differ in how they execute a module. Wasmtime, the Bytecode
Alliance's runtime, compiles modules with Cranelift and implements the
WASI system interface \cite{wasmtime,wasi}. WAMR (WebAssembly Micro
Runtime) is a small C runtime for embedding, with a classic and a fast
interpreter and optional ahead-of-time and just-in-time modes
\cite{wamr}. \paglets{}/cpp embeds WAMR in its worker processes and runs
the fast interpreter by default (the classic interpreter is a build
option); neither the ahead-of-time nor the just-in-time modes are built
(\code{cpp/cmake/PagletsWamr.cmake}). The interpreter costs speed
(Section~\ref{sec:evaluation}). It is the portable default: WAMR's
ahead-of-time mode needs an LLVM-based compiler (\code{wamrc}) and has
not been measured, and the design keeps compiled forms, where they exist,
in a per-host cache that is never sent between hosts \cite{pagletsCpp}.
The image is independent of the execution mode either way: it holds the
linear memory and the globals, and no code. Of WASI a \paglet{} may
import only a subset, listed in Appendix~\ref{app:abi}; everything else
it reaches through the \code{paglets} import module.

The idea of treating an instance's memory as a snapshot is not new. Wizer
runs a module's initialization ahead of time and writes the resulting
memory and globals into the data segments of a new module, so that the
module starts initialized \cite{wizer}. It uses the same observation for
a different purpose; the initialization it snapshots may not call
imported functions (WASI calls only when explicitly allowed), and the
module may not import globals, tables or memories. Research on
checkpointing and migrating running Wasm programs attacks the harder
case: Tamura et al. instrument ahead-of-time compiled code and record
stack maps, so that at designated program points the machine state can
be turned back into Wasm state (operand stack and locals) and restored on
another instruction set \cite{tamura2025crossisa}. \paglets{} needs none
of this because it moves at quiescent points only; it trades the ability
to move mid-computation for an image that is a plain copy of memory. Long
computations are therefore written in steps, as the examples in
Sections~\ref{sec:zombies} and \ref{sec:evaluation} show; the host
enforces this with a time budget per handler call (5~s by default), after
which the \paglet{} fails.

\subsection{The Python edition and what the C++ edition changes}\label{sec:background-python}

The first edition of \paglets{} is a Python package
\cite{pagletsSoftware,klukas2026pagletsNote}. Its mobile state is a
dataclass the author declares, and a move carries that state with the
names of the paglet and state classes, which every host must be able to
import; code is never uploaded. Each active paglet runs in a child
process of the host, and the hosts of a mesh trust each other and may
share an API key.

The C++ edition keeps the concepts (hosts, paglets, messages, a mesh,
services, compute slots) and none of the formats; the two editions do not
exchange paglets \cite{pagletsCppDocs}. It changes five things.

\begin{enumerate}
\item \emph{Modules are the unit of code.} A \paglet{}'s code is a
  WebAssembly module, named by its SHA-256 hash. A host that lacks a
  module fetches it from the source host or any peer and verifies the
  hash; hosts do not need the same code installed
  (Section~\ref{sec:movement}).
\item \emph{State is the memory image.} There is no state class. The
  author's objects live in linear memory, and a move sends the pages
  that changed since the destination last saw them, compressed.
\item \emph{Paglets run sandboxed.} Instances live in worker processes
  that the operating system confines (seccomp on Linux, a sandbox
  profile on macOS, a job object with a restricted token on Windows).
  A \paglet{} cannot open a file or a socket; it reaches host resources
  through system paglets, within the mesh policy
  (Section~\ref{sec:mesh}).
\item \emph{A signed ledger replaces the API key.} Hosts, admins, owners
  and signers have Ed25519 keys. Membership, module trust, policy rules
  and grants are signed records, spread by gossip; hosts talk inside
  mutually authenticated, end-to-end encrypted channels
  (Section~\ref{sec:mesh}).
\item \emph{Fault tolerance with a checked protocol.} Zombie paglets
  replicate a job on every host of the mesh that offers the
  \code{zombies} service and hand it over when the main replica
  disappears. The protocol is a C++ state machine that the guest SDK runs
  and that an explicit-state model checker explores for small worlds
  (Section~\ref{sec:zombies}).
\end{enumerate}

Table~\ref{tab:background-editions} places the three systems next to each
other by concept; Section~\ref{sec:discussion-comparison} does it again
by what each one pays, with the limits that the evaluation shows.

\begin{table}[htbp]
\centering\small
\caption{Aglets and the two \paglets{} editions, by concept.}
\label{tab:background-editions}
\begin{tabular}{p{0.16\linewidth}p{0.23\linewidth}p{0.23\linewidth}p{0.23\linewidth}}
\toprule
 & Aglets & \paglets{}, Python & \paglets{}/cpp \\
\midrule
Mobile unit & Java object & dataclass state plus class names & Wasm instance: module hash plus memory image \\
Execution & \code{run()} thread, restarted after a move & a \code{run()} hook in the child process, plus handlers, serial per paglet & handlers, one at a time; no thread \\
Reaching others & proxies & proxies & endpoint capabilities only \\
Failure & --- & an active paglet is lost with its host; an inactive one survives as persisted dataclass state & a trap ends the paglet; after a host restart it resumes from its last image (default policy); zombie replicas \\
\bottomrule
\end{tabular}
\end{table}
